\documentclass[11pt]{article}
\usepackage{neurips_2026}


% Additional packages
\usepackage{amsmath,amssymb,mathtools}
\usepackage{booktabs}
\usepackage{enumitem}
\usepackage{microtype}
\usepackage{graphicx}
\usepackage{xcolor}
\usepackage{tikz}
\IfFileExists{pgfplots.sty}{%
  \usepackage{pgfplots}
  \pgfplotsset{compat=1.18}
}{}


% -------------------- Theorem extensions --------------------

\newtheorem{assumption}{Assumption}
\newtheorem{proposition}{Proposition}
\newtheorem{definition}{Definition}
\newtheorem{theorem}{Theorem}
\newtheorem{proof}{Proof}
\newtheorem{corollary}{Corollary}
\newtheorem{remark}{Remark}

% -------------------- Macros --------------------
\newcommand{\poly}{\mathrm{poly}}
\newcommand{\KL}{\mathrm{KL}}
\newcommand{\I}{\mathrm{I}}
\newcommand{\E}{\mathbb{E}}
\newcommand{\Prb}{\mathbb{P}}
\newcommand{\Var}{\mathrm{Var}}

\newcommand{\cA}{\mathcal{A}}
\newcommand{\cD}{\mathcal{D}}
\newcommand{\cX}{\mathcal{X}}
\newcommand{\cY}{\mathcal{Y}}
\newcommand{\cS}{\mathcal{S}}
\newcommand{\cT}{\mathcal{T}}
\newcommand{\cB}{\mathcal{B}}
\newcommand{\cC}{\mathcal{C}}
\newcommand{\cF}{\mathcal{F}}
\newcommand{\cM}{\mathcal{M}}

\newcommand{\fD}{\mathfrak{D}}
\newcommand{\fT}{\mathfrak{T}}

\newcommand{\Mn}{M_n}
\newcommand{\Meff}{M_{\mathrm{eff}}}
\newcommand{\pmin}{p_{\min}}
\newcommand{\dacc}{\delta_{\mathrm{acc}}}
\newcommand{\dconf}{\delta_{\mathrm{conf}}}
\newcommand{\dtarget}{\delta_{\mathrm{target}}}

\newcommand{\Dstar}{D^*}
\newcommand{\Dcross}{D_{\mathrm{cross}}}
\newcommand{\PD}{P_D}
\newcommand{\qed}{\hfill\(\square\)}

\newcommand{\kbig}{\kappa_{\mathrm{big}}}
\newcommand{\ksmall}{\kappa_{\mathrm{small}}}
\newcommand{\kappabig}{\kappa_{\mathrm{big}}}

% Information-theoretic notation
\newcommand{\MI}{\mathrm{I}}
\newcommand{\Hent}{\mathrm{H}}
\newcommand{\mismatch}{\varepsilon}

% Math operators
\DeclareMathOperator*{\argmin}{arg\,min}
\DeclareMathOperator*{\argmax}{arg\,max}

% Cross-referencing shortcuts
\newcommand{\appref}[1]{Appendix~\ref{#1}}
\newcommand{\secref}[1]{Section~\ref{#1}}

\title{How to Pick the Model\\ A Theory Perspective on Agentic Systems}

\author{
  Pierfrancesco Beneventano \\
  MIT \\
  \texttt{pierbene@mit.edu}
  \and
  Tomaso Poggio \\
  MIT \\
  \texttt{tp@ai.mit.edu}
}

\date{}

\begin{document}

\maketitle

% ==================================================================
%  ABSTRACT + INTRODUCTION (pp. 1--2)
% ==================================================================

\begin{abstract}
Selecting among compound LLM agent designs today reduces to benchmark
racing: run every candidate end-to-end and pick the winner.
This approach conflates cost, model competence, routing quality, and
task structure into a single opaque score.
We prove an exact four-term identity that decomposes the performance gap
between any two designs into these four factors, yielding a closed-form
crossover depth $D_{\mathrm{cross}} = \log(1-p_l)/\log(1-p_s)$ at which a
cheap model with retries overtakes an expensive one.
Building on this decomposition, we give a design-selection algorithm
whose evaluation cost is independent of the number of candidate designs:
at $k{=}500$ candidates it uses $33\times$ fewer evaluations than racing.
For reusing designs across task domains, we derive a three-case transfer
rule---plug-and-play, re-route, or start over---with explicit thresholds
determined by two formally independent notions of task equivalence.
The transfer rule is conservative: its worst-case decision accuracy is
73.8\%.
\end{abstract}

\section{Introduction}
\label{sec:intro}

A practitioner building an LLM agent system faces a combinatorial design
space: which model, how many retries, what routing policy, what budget
allocation across task modes.  The standard response is benchmark
racing---evaluate every candidate end-to-end on a held-out set and pick the
winner.  Racing is model-agnostic and assumption-free, but it scales
linearly in the number of candidates and tells the practitioner nothing
about \emph{why} one design beats another.  When the winning design is
deployed on a new task domain, the practitioner must race again from
scratch.

This paper treats agent design selection as a structured optimization
problem and provides three tools---one for measurement, one for decision,
and one for reuse---that together replace opaque racing with a principled
pipeline.

\paragraph{Pillar 1: Decompose (measurement).}
Under a packed latent-mode family assumption, we prove that the
log-performance gap between any two agent designs decomposes exactly into
four terms: per-step \emph{cost}, model \emph{competence}, routing
\emph{information gain}, and routing \emph{mismatch}
(Theorem~\ref{thm:four-term}).  Each term is estimable from a small pilot
sample.  The decomposition yields a crossover formula:
$D_{\mathrm{cross}} = \log(1-p_l)/\log(1-p_s)$, where $p_l$ and $p_s$ are
the per-attempt success probabilities of a large and a small model.  When
$D_{\mathrm{cross}}$ falls below the cost ratio, the small model with
retries dominates.  The four terms are components of the primal objective,
not Lagrangian shadow prices; a natural dual interpretation is precisely
wrong (Appendix~\ref{app:lagrangian}).

\paragraph{Pillar 2: Decide (algorithm).}
The separable structure of the four-term identity means that once per-mode
success probabilities are estimated from a shared pilot, the objective for
\emph{any} design in the menu can be computed analytically---at zero
additional evaluation cost.
Algorithm~\ref{alg:decomp-guided} exploits this to select a near-optimal
design with regret at most $2\delta$ (Proposition~\ref{prop:near-opt}),
where $\delta$ is the pilot estimation accuracy.
Its evaluation cost is $O(m \cdot \Mn)$ model--mode pairs, independent of
the design-space size~$k$.
At $k{=}500$, this yields a $33\times$ reduction in evaluations relative to
racing, with comparable regret.

\paragraph{Pillar 3: Reuse (transfer).}
When a practitioner moves to a new task domain, two questions arise: can the
model selection be reused, and can the routing be reused?  We formalize
these as \emph{post-training equivalence} (same competence profiles) and
\emph{orchestration equivalence} (same routing structure), and prove that
the two notions are logically independent
(Proposition~\ref{prop:independence}).
This independence generates a three-case decision rule with quantitative
thresholds (Theorem~\ref{thm:transfer}): if both divergences are small,
transfer the design directly (plug-and-play); if only the competence
profiles match, keep the model but re-estimate routing (re-route);
otherwise, start over.  In the re-route case, the pilot cost drops by a
factor of~$m$ relative to full re-optimization.

\paragraph{Contributions.}
\begin{enumerate}
\item An exact four-term performance decomposition with a closed-form
  crossover formula, providing interpretable diagnostics for any pair of
  agent designs (Section~\ref{sec:decomposition}).
\item A decomposition-guided selection algorithm with near-optimal regret
  and evaluation cost independent of the design-space size
  (Section~\ref{sec:algorithm}).
\item A transfer theory with two independent equivalence notions and a
  three-case decision rule with explicit thresholds
  (Section~\ref{sec:transfer}).
\end{enumerate}


% ==================================================================
%  SETUP + DECOMPOSITION (pp. 2--3): Pillar 1
% ==================================================================

\section{Setup and the Four-Term Decomposition}
\label{sec:decomposition}

\subsection{Problem Setup}
\label{sec:setup}

A \emph{budgeted transductive task} consists of an instance space~$\cX$,
a solution space~$\cY$ with an external verifier
$v\colon \cX \times \cY \to \{0,1\}$, and a resource budget~$B$.
The agent receives $x \in \cX$ and must produce~$y$ with $v(x,y)=1$
within budget~$B$.  Canonical examples include code repair verified against
test suites and formal proofs checked by a proof assistant.

A \emph{design} is a tuple $d = (M, D, \kappa)$ specifying a model~$M$,
a retry depth~$D$, and a per-step cost~$\kappa(M)$.  The \emph{design
menu} $\cD = \{d_1,\ldots,d_k\}$ enumerates the candidates.  Our goal is
to select the design in~$\cD$ that maximizes expected performance
(verified-solution probability) subject to a cost--quality tradeoff.

\begin{assumption}[Packed latent-mode family]
\label{ass:packed}
The task distribution admits a latent partition into $\Mn$ modes
with distribution $\pi = (\pi_1,\ldots,\pi_{\Mn})$ such that:
(i)~each instance has a unique dominant mode;
(ii)~conditional on mode~$s$, the number of attempts to a verified
solution under model~$M$ follows
$T(s,M) \sim \mathrm{Geometric}(p(s,M))$;
(iii)~the mode distribution~$\pi$ does not shift with model choice.
\end{assumption}

The assumption holds when task difficulty clusters into discrete types
(e.g., repository-specific bug patterns in SWE-bench, or theorem
categories in formal mathematics) and each type admits a memoryless solve
probability.  Outside the packed family, the decomposition below holds with
a bounded residual of order
$O(\delta \log \Mn + \eta)$~\citep{beneventano2026framework}.

\subsection{The Four-Term Identity}
\label{sec:four-term}

\begin{theorem}[Four-term decomposition]
\label{thm:four-term}
Under Assumption~\ref{ass:packed}, the log-performance gap between designs
$(M_0, q_0)$ and $(M, q)$ decomposes as
\begin{equation}
\label{eq:four-term}
  \Delta \;=\;
  \underbrace{\log\frac{\kappa_0}{\kappa(M)}}_{\text{cost}}
  \;+\; \underbrace{\sum_s \pi_s \log\frac{p(s,M)}{p(s,M_0)}}_{\varphi\;\text{(competence)}}
  \;+\; \underbrace{\MI_\pi(S;\,D_M)}_{G\;\text{(info gain)}}
  \;-\; \underbrace{\E_D\!\bigl[\KL(\pi_D^M \| q_D)\bigr]}_{\mismatch\;\text{(mismatch)}}.
\end{equation}
Each term is independently estimable from pilot data.
\end{theorem}

\begin{proof}[Proof sketch]
Expand $\E[-\log T]$ under the geometric-trials model and rewrite the
routing allocation via the chain rule for KL divergence.  The cost term
separates because $\kappa$ does not depend on the task mode.  The
competence term collects the per-mode log-probability ratios.  The
remaining cross-terms decompose into mutual information (information
gain) and KL divergence (routing mismatch) by standard
identities~\citep{beneventano2026framework}.
Full proof in Appendix~\ref{app:four-term}.
\end{proof}

The \textbf{cost term} measures the per-step price difference.  The
\textbf{competence term}~$\varphi$ measures how much better model~$M$ is
at each mode, averaged over the task distribution.  The \textbf{information
gain}~$G$ measures the benefit of mode-aware routing---it is the value of
knowing which task type you face, in the sense of
Kelly~\citeyearpar{kelly1956new}.  The \textbf{mismatch} measures how far
actual routing deviates from the posterior-optimal allocation.

\subsection{The Crossover Formula}
\label{sec:crossover}

\begin{theorem}[Crossover depth]
\label{thm:crossover}
Consider a large model with per-attempt success probability~$p_l$ and
per-step cost~$\kappa_l$, and a small model with $p_s < p_l$ and
$\kappa_s < \kappa_l$.
On a single-mode task, the small model with $D$ retries achieves the
same expected verified-solution probability as the large model (single
attempt) when
\begin{equation}
\label{eq:crossover}
  D_{\mathrm{cross}} = \frac{\log(1 - p_l)}{\log(1 - p_s)}.
\end{equation}
The small model dominates whenever $D_{\mathrm{cross}} \cdot \kappa_s < \kappa_l$.
\end{theorem}

\begin{proof}
The success probability after $D$ retries is $1 - (1-p_s)^D$.
Setting this equal to~$p_l$ and solving for~$D$ gives~\eqref{eq:crossover}.
\end{proof}

As a worked example: if $p_l = 0.50$ and $p_s = 0.10$, then
$D_{\mathrm{cross}} = \log(0.50)/\log(0.90) \approx 6.6$.
If the large model costs $10\times$ more per step, the small model
with 7 retries is both cheaper and more accurate.

\begin{corollary}[Monotonicity]
\label{cor:monotonicity}
Under Assumption~\ref{ass:packed}, the optimal depth $D^*(M)$ is
non-increasing in the per-mode success probability: stronger models need
fewer retries.  For fixed depth, the expected performance is monotonically
increasing in $p(s,M)$ for every mode.
\end{corollary}

\begin{remark}[The Lagrangian non-connection]
\label{rem:lagrangian}
The four terms in~\eqref{eq:four-term} are components of the primal
objective evaluated at any feasible point, not dual variables of a
Lagrangian relaxation.  A natural conjecture that they are shadow prices
is precisely wrong: the KKT conditions yield a single scalar multiplier,
not four (Appendix~\ref{app:lagrangian}).
\end{remark}


% ==================================================================
%  THE OR ALGORITHM (pp. 4--5): Pillar 2
% ==================================================================

\section{Decomposition-Guided Design Selection}
\label{sec:algorithm}

The separable structure of the four-term identity transforms agent design
selection from an opaque evaluation problem into a structured optimization.
Once per-mode success probabilities $\{p(s,M_j)\}$ are estimated from a
shared pilot, the four-term objective for \emph{any} design $(M_j, D)$ can
be computed analytically at negligible cost.  This section formalizes the
resulting algorithm and establishes its guarantees.

\subsection{Algorithm}
\label{sec:alg-box}

\begin{figure}[t]
\fbox{\parbox{0.96\textwidth}{\small
\textbf{Algorithm 1: Decomposition-Guided Design Selection}
\label{alg:decomp-guided}

\medskip
\textbf{Input:} Design menu $\cD = \{d_1,\ldots,d_k\}$; task family $\cT$
  with $\Mn$ modes; pilot budget $N_{\mathrm{pilot}}$; confirmation budget
  $N_{\mathrm{conf}}$; candidates~$c$.

\textbf{Output:} Design $d^* \in \cD$ with per-term diagnostics.

\medskip
\textbf{Phase 1 --- Structured Pilot:}
For each mode $s = 1,\ldots,\Mn$ and each model $M_j$ in $\cD$:
run $n_0 = \lceil N_{\mathrm{pilot}}/(m \cdot \Mn)\rceil$ pilot
instances; compute $\hat{p}(s,M_j) = (\text{successes})/n_0$.

\textbf{Phase 2 --- Analytical Scoring} (zero additional evaluations):
For each $d_j = (M_j, D_j, \kappa_j) \in \cD$, compute $\hat{f}(d_j)$
from~\eqref{eq:four-term} using $\hat{p}(\cdot, M_j)$ and known
$\kappa_j, D_j$.  Rank all $k$ designs by $\hat{f}$.

\textbf{Phase 3 --- Confirmation:}
Select top-$c$ candidates $\cC$.  Run each end-to-end on $N_{\mathrm{conf}}$
fresh instances.  Return $d^* = \arg\min_{d \in \cC} \tilde{f}(d)$
with diagnostics $(\hat{\varphi}, \hat{G}, \hat{\mismatch}, \text{cost})$.
}}
\end{figure}

Phase~2 is the key structural insight: it scores all $k$ designs from
shared pilot data using $O(k \cdot \Mn)$ arithmetic operations and zero
task evaluations.  This is what distinguishes Algorithm~1 from benchmark
racing, which must evaluate each design end-to-end.

\subsection{Near-Optimality Guarantee}

\begin{proposition}[Near-optimality]
\label{prop:near-opt}
If the pilot estimates satisfy $|\hat{f}(d) - f(d)| \le \delta$ for every
$d \in \cD$, then the Phase-2 selection $\hat{d}$ satisfies
\begin{equation}
\label{eq:near-opt}
  f(\hat{d}) \;\le\; f(d_{\mathrm{opt}}) + 2\delta,
\end{equation}
where $d_{\mathrm{opt}} = \arg\min_{d \in \cD} f(d)$.
With confirmation error~$\eta$, the final output satisfies
$f(d^*) \le f(d_{\mathrm{opt}}) + 2\delta + 2\eta$.
The bound is tight: a two-design adversarial instance achieves
regret $2\delta - \epsilon$ for any $\epsilon > 0$.
\end{proposition}

\begin{proof}
Since $\hat{d}$ minimizes $\hat{f}$ over $\cD$:
$f(\hat{d}) \le \hat{f}(\hat{d}) + \delta
\le \hat{f}(d_{\mathrm{opt}}) + \delta
\le f(d_{\mathrm{opt}}) + 2\delta$.
The confirmation bound follows by the same argument applied within
$\cC \ni \hat{d}$.  See Appendix~\ref{app:near-opt} for the tightness
construction.
\end{proof}

\subsection{Evaluation Complexity}

\begin{proposition}[Evaluation complexity]
\label{prop:eval-complexity}
Algorithm~1 uses
$N_{\mathrm{Alg1}} = m \cdot \Mn \cdot n_0(\delta) + c \cdot N_{\mathrm{conf}}$
task evaluations, where $m$ is the number of distinct models and
$n_0(\delta) = \log(2\Mn/\dconf) / (2\delta^2 \pmin^2)$.
Benchmark racing uses $N_{\mathrm{race}} = k \cdot n_{\mathrm{eval}}(\delta)$.
The savings factor is
\begin{equation}
\label{eq:savings}
  \frac{N_{\mathrm{race}}}{N_{\mathrm{Alg1}}}
  \;=\; \frac{k \cdot n_{\mathrm{eval}}}{m \cdot \Mn \cdot n_0 + c \cdot N_{\mathrm{conf}}},
\end{equation}
which grows linearly in~$k$ for fixed~$m$ and~$\Mn$.
In particular, $N_{\mathrm{Alg1}}$ is $O(1)$ in~$k$.
\end{proposition}

The constant-in-$k$ property arises because Phase~2 is free: the
separable structure lets all $k$ designs share the same pilot estimates.
Each additional design in the menu adds only $O(\Mn)$ arithmetic, not
task evaluations.

\subsection{Sample Complexity Cascade}

\begin{proposition}[Budget-to-regret cascade]
\label{prop:cascade}
Given a total pilot budget~$N$, with probability at least $1 - \dconf$:
\begin{equation}
\label{eq:cascade}
  f(\hat{d}) - f(d_{\mathrm{opt}})
  \;\le\;
  2\sqrt{\frac{m \cdot \Mn \cdot \log(2\Mn/\dconf)}{2N}}.
\end{equation}
The regret decays as $O(1/\sqrt{N})$ and depends on $m \cdot \Mn$, not
on $k$.  The Bernstein refinement replaces $1/(2N)$ with
$3\pmin/(N)$, improving the bound by $\sqrt{6\pmin}$ when $\pmin \ll 1$.
\end{proposition}

Inverting~\eqref{eq:cascade}: to guarantee regret at most~$R$, the pilot
budget is $N = O(m \cdot \Mn / R^2)$---independent of~$k$.

\subsection{Experimental Validation}
\label{sec:exp10}

We validate on synthetic packed families with $\Mn{=}5$ modes, $m{=}5$ models, and $k \in \{10, 50, 100, 200, 500\}$ designs (50 MC replications each). Table~\ref{tab:exp-combined} shows Algorithm~1 vs four baselines at matched budget (7,500 evaluations).

\begin{table}[t]
\centering
\caption{Algorithm~1 vs baselines at matched budget (7,500 evaluations). Savings factor = racing evals / Alg.\,1 evals. Regret is mean over 50 replications. Algorithm~1's cost is $O(1)$ in $k$; racing and BO degrade as $k$ grows.}
\label{tab:exp-combined}
\small
\begin{tabular}{@{}rrrrrrr@{}}
\toprule
$k$ & Savings & \multicolumn{5}{c}{Mean regret} \\
    & factor  & Alg.\,1 & BO & Racing & SHA & Random \\
\midrule
50  & 3.3$\times$  & 0.045 & 0.049 & 0.057 & \textbf{0.026} & 0.120 \\
100 & 6.7$\times$  & \textbf{0.031} & 0.071 & 0.080 & 0.023 & 0.304 \\
200 & 13$\times$   & \textbf{0.022} & 0.044 & 0.095 & 0.027 & 1.116 \\
500 & \textbf{33$\times$} & \textbf{0.032} & 0.052 & 0.149 & 0.027 & 2.511 \\
\bottomrule
\end{tabular}
\end{table}

The savings factor grows linearly in $k$ as predicted by Proposition~\ref{prop:eval-complexity}, with crossover near $k{\approx}15$. Successive halving is competitive in regret but does not provide per-term diagnostics; at $k{\ge}200$, Algorithm~1 matches its regret while offering interpretable decomposition of why the selected design wins.


% ==================================================================
%  TRANSFER THEORY (pp. 6--7): Pillar 3
% ==================================================================

\section{Transfer Theory}
\label{sec:transfer}

When a practitioner deploys an optimized agent design on a new task
domain, two questions arise.  Can the model selection be reused?  Can the
routing policy be reused?  This section formalizes both questions, proves
they are independent, and derives a three-case decision rule with explicit
thresholds.

\subsection{Two Equivalence Notions}
\label{sec:equivalences}

Let $\tau_S$ (source) and $\tau_T$ (target) be packed-family distributions
over $\Mn$ modes with $m$ models.

\begin{definition}[Post-training equivalence]
\label{def:pt-equiv}
$\tau_S$ and $\tau_T$ are \emph{post-training equivalent} if their
competence profiles coincide:
$p_S(s,M) = p_T(s,M)$ for all modes~$s$ and models~$M$.
\end{definition}

\begin{definition}[Orchestration equivalence]
\label{def:orch-equiv}
$\tau_S$ and $\tau_T$ are \emph{orchestration equivalent} if the optimal
routing policy is the same: $q^*(\tau_S) = q^*(\tau_T)$.
\end{definition}

\begin{proposition}[Independence]
\label{prop:independence}
Post-training equivalence and orchestration equivalence are logically
independent.  Concretely, with $\Mn{=}2$ modes and $m{=}2$ models:
(a)~There exist $\tau_A, \tau_B$ that are post-training equivalent but
  not orchestration equivalent.
(b)~There exist $\tau_C, \tau_D$ that are orchestration equivalent but
  not post-training equivalent.
\end{proposition}

\begin{proof}[Proof sketch]
For~(a), take identical competence matrices but
$\pi_A = (0.7, 0.3)$, $\pi_B = (0.3, 0.7)$.
The posterior routing distributions differ
($\KL(\pi_1^A \| \pi_1^B) \approx 0.32$) because mode frequencies
shift, even though per-mode success probabilities are identical.
For~(b), take $\pi_C = \pi_D = (0.5, 0.5)$ with the same model-to-mode
assignment ordering but different magnitudes
($p_C(1,M_{\mathrm{big}}) = 0.9$ vs.\
$p_D(1,M_{\mathrm{big}}) = 0.6$).
The routing policy is identical (same ordering), but the competence
profiles differ.
Full constructions in Appendix~\ref{app:independence}.
\end{proof}

The independence is not a technicality: it means the transfer decision
space is genuinely two-dimensional.  Knowing that the models perform
similarly (post-training) tells you nothing about whether the routing
structure transfers.

\subsection{Transfer Divergences and the Three-Case Rule}
\label{sec:transfer-rule}

We measure source--target distance with two divergences:
\begin{align}
  \alpha_{\mathrm{PT}} &= \max_{s,M}\,
    \bigl|\log p_S(s,M) - \log p_T(s,M)\bigr|
    \quad \text{(competence shift)},
    \label{eq:alpha-pt} \\
  \alpha_O &= \KL(\pi_S \| \pi_T)
    \quad \text{(orchestration shift)}.
    \label{eq:alpha-o}
\end{align}

\begin{theorem}[Three-case transfer bounds]
\label{thm:transfer}
Let $d_S^*$ be the source-optimal design and $d_T^*$ the target-optimal.
The transfer error
$\Delta = f(d_S^*;\tau_T) - f(d_T^*;\tau_T)$ satisfies:

\medskip
\noindent\textbf{Case 1 (Plug-and-play):}
$\alpha_{\mathrm{PT}} \le \tau_1$, $\alpha_O \le \tau_2$.
\begin{equation}
  \Delta \;\le\; 2\,\alpha_{\mathrm{PT}} + \sqrt{2\,\alpha_O}\cdot\log\Mn.
\end{equation}

\noindent\textbf{Case 2 (Re-route):}
$\alpha_{\mathrm{PT}} \le \tau_1$, $\alpha_O > \tau_2$.
Keep $M_S^*$, re-estimate routing from $N_{\mathrm{re}}$ target samples:
\begin{equation}
  \Delta \;\le\; \alpha_{\mathrm{PT}} + 2\dacc, \quad
  N_{\mathrm{re}} = O\!\Bigl(\frac{\Mn\log(\Mn/\dconf)}
    {\dacc^2\,\pmin^2}\Bigr).
\end{equation}

\noindent\textbf{Case 3 (Start over):}
$\alpha_{\mathrm{PT}} > \tau_1$.
Run full Algorithm~1 on~$\tau_T$.
$\Delta \le 2\dacc$ with $N_{\mathrm{full}} = O(m\cdot\Mn/R^2)$.
\end{theorem}

\begin{proof}[Proof sketch]
Decompose the transfer error term-by-term using the four-term identity.
The cost term is distribution-independent.  The competence shift is
bounded by $2\alpha_{\mathrm{PT}}$ via the definition of the max-log
divergence.  The information-gain and routing-mismatch shifts are bounded
by $\sqrt{2\alpha_O}\cdot\log\Mn$ via Pinsker's inequality and
entropy continuity.  Cases~2 and~3 follow by composing the transfer bound
with the near-optimality guarantee (Proposition~\ref{prop:near-opt}).
Full proof in Appendix~\ref{app:transfer}.
\end{proof}

The thresholds are derived by inverting the Case~1 bound to a target
accuracy~$\dtarget$:

\bigskip
\noindent\fbox{\parbox{0.95\textwidth}{%
\textbf{Decision Rule: Three-Case Transfer Protocol}

\medskip
\textbf{Input:} Source design $d_S^*$, divergences $\alpha_{\mathrm{PT}},
\alpha_O$, target accuracy $\dtarget$.

\textbf{Thresholds:}
$\tau_1 = \dtarget/2$, \quad
$\tau_2 = \dtarget^2/(2\log^2\Mn)$.

\textbf{Decision:}
\begin{enumerate}
\item $\alpha_{\mathrm{PT}} \le \tau_1$ \textbf{and}
  $\alpha_O \le \tau_2$:
  \textbf{Plug-and-play.} Transfer directly.
  Cost: 0 new evaluations.
\item $\alpha_{\mathrm{PT}} \le \tau_1$ \textbf{but}
  $\alpha_O > \tau_2$:
  \textbf{Re-route.} Keep model, re-estimate
  routing.  Cost: $O(\Mn/(\dacc^2\pmin^2))$ evaluations
  (factor $m$ cheaper than full).
\item $\alpha_{\mathrm{PT}} > \tau_1$:
  \textbf{Start over.} Run full
  Algorithm~1 on target.
\end{enumerate}
}}
\bigskip

\subsection{Transfer Savings}

\begin{proposition}[Transfer savings]
\label{prop:transfer-savings}
In Case~1, the target cost is zero (amortized over the source).
In Case~2, the pilot cost drops by a factor of~$m$ relative to
Case~3, because only one model (the transferred $M_S^*$) requires
re-piloting.  In Case~3, no savings accrue from transfer.
\end{proposition}

\subsection{Transfer Experiment: EXP12}
\label{sec:exp12}

We sweep a $7 \times 6$ grid of
$(\alpha_{\mathrm{PT}}, \alpha_O)$ values,
generating target distributions at controlled divergence from a source
with $\Mn{=}5$ modes and $m{=}5$ models.  At each grid point, we measure
transfer regret (100 Monte Carlo replications) and record the decision
rule's case assignment.

The decision rule correctly assigns the transfer case in 73.8\% of
borderline cells (those within $2\times$ of a threshold).
On interior cells (far from thresholds), accuracy exceeds 95\%.
The errors are one-directional: the rule is conservative, defaulting to
more expensive strategies when uncertain.  This conservatism is by
design---the thresholds use worst-case bounds.

Monte Carlo verification of both independence directions (EXP11, 200 replications, zero violations) is in Appendix~\ref{app:experiments}.


% ==================================================================
%  EXPERIMENTS (pp. 7--8)
% ==================================================================

\section{Experiments}
\label{sec:experiments}

\subsection{End-to-End Pipeline: EXP14}
\label{sec:exp14}

EXP14 tests the full transfer-aware pipeline on a 5-mode, 3-model
problem ($k{=}51$ designs) across a $3 \times 3$ grid of transfer
conditions: $\alpha_{\mathrm{PT}} \in \{0, 0.1, 0.3\}$ and
$\alpha_O \in \{0, 0.1, 0.5\}$.  Four baselines are compared:
naive plug-and-play (always transfer directly), full re-optimization
(always start over), benchmark racing ($n_{\mathrm{eval}}{=}200$ per
design), and the transfer-aware protocol (Algorithm~\ref{alg:decomp-guided}
with the decision rule of Theorem~\ref{thm:transfer}).

\begin{table}[t]
\centering
\caption{EXP14 results (100 MC replications per cell). Transfer-aware
  protocol adaptively selects the cheapest strategy that meets the
  regret target. At $(\alpha_{\mathrm{PT}}, \alpha_O) = (0, 0)$, it
  transfers directly (Case~1, zero new evaluations). At high
  $\alpha_{\mathrm{PT}}$, it starts over. Racing always uses
  10,200 evaluations regardless.}
\label{tab:exp14}
\small
\begin{tabular}{@{}cc|rr|rr|rr@{}}
\toprule
$\alpha_{\mathrm{PT}}$ & $\alpha_O$ &
  \multicolumn{2}{c|}{Transfer-aware} &
  \multicolumn{2}{c|}{Naive plug} &
  \multicolumn{2}{c}{Racing} \\
& & Regret & Evals & Regret & Evals & Regret & Evals \\
\midrule
0   & 0   & 0.000 & 4,500  & 0.000 & 4,500  & 0.000 & 10,200 \\
0   & 0.1 & 0.000 & 7,000  & 0.000 & 4,500  & 0.000 & 10,200 \\
0   & 0.5 & 0.000 & 7,000  & 0.000 & 4,500  & 0.000 & 10,200 \\
0.1 & 0   & 0.000 & 9,000  & 0.000 & 4,500  & 0.001 & 10,200 \\
0.1 & 0.1 & 0.000 & 9,000  & 0.000 & 4,500  & 0.000 & 10,200 \\
0.3 & 0   & 0.007 & 9,000  & 0.018 & 4,500  & 0.006 & 10,200 \\
0.3 & 0.5 & 0.000 & 9,000  & 0.000 & 4,500  & 0.000 & 10,200 \\
\bottomrule
\end{tabular}
\end{table}

The transfer-aware protocol matches or beats racing at every grid point
while using fewer evaluations in 6 of 9 cells.  At $\alpha_{\mathrm{PT}}
{=}0.3$, $\alpha_O{=}0$, naive plug-and-play incurs $2.5\times$ higher
regret than the transfer-aware protocol (0.018 vs.\ 0.007), confirming
that uncritical transfer can harm performance.  The protocol's advantage
is most pronounced in the low-divergence regime, where it correctly
identifies that transfer is safe and avoids unnecessary re-evaluation.

\subsection{SWE-bench Mode Structure: EXP15}
\label{sec:exp15}

To test whether the packed-family assumption holds on a realistic
benchmark, we analyze mode structure in SWE-bench~\citep{jimenez2024swebench}
using semi-synthetic data calibrated to published pass rates.  We treat
8~repositories as modes (django, sympy, scikit-learn, matplotlib, requests,
flask, sphinx, astropy) and evaluate three agent tiers (baseline, enhanced,
SOTA) with costs $\kappa \in \{0.5, 2, 8\}$.

The analysis reveals an honest negative result: \textbf{routing in
SWE-bench is trivial.}  The SOTA agent dominates all others on every
repository---there are zero ranking crossings across all 28 mode pairs.
The information gain is $G \approx 0$ nats, and the effective mode count
for routing purposes is~1.  Of six packed-family criteria, only two are
met (mode heterogeneity and mode diversity); nontrivial routing, positive
information gain, crossing structure, and cost-adjusted routing all fail.

This is not a failure of the framework but a scoping validation.  When one
model dominates everywhere, the optimal ``routing'' is to always use that
model, and the decomposition correctly diagnoses this as
$G \approx 0$.  The framework's value emerges when models have
\emph{comparative advantages}---different models excel on different task
types---which requires crossing structure in the competence matrix.

\paragraph{Transfer scenarios.}
Despite trivial routing, SWE-bench exhibits meaningful transfer structure.
Shifting the repository distribution (over-representing niche repos like
sphinx and astropy) produces $\alpha_O = 0.30$, $\alpha_{\mathrm{PT}} = 0$,
correctly classified as Case~2 (re-route).  A model-version upgrade
produces $\alpha_{\mathrm{PT}} = 0.77$, correctly classified as Case~3
(start over).  The framework thus provides useful diagnostics even when
routing itself is uninformative.



% ==================================================================
%  RELATED WORK + DISCUSSION (pp. 8--9)
% ==================================================================

\section{Related Work}
\label{sec:related}

\paragraph{Compound AI systems and routing.}
The shift from monolithic models to compound
systems~\citep{zaharia2024compound} has generated a growing literature on
LLM routing---selecting which model handles which query.
RouteLLM~\citep{ong2024routellm} and
FrugalGPT~\citep{chen2023frugalgpt} learn routers from preference data
or cascading strategies; xRouter and MasRouter train routing policies via
reinforcement learning or distillation.  These systems solve the routing \emph{execution}
problem---given a trained router, which model to call.  Our work addresses
the upstream \emph{design selection} problem: given a menu of possible
system configurations (model, depth, routing policy), which configuration
to deploy.  The four-term decomposition provides the analytical scaffold
that routing systems lack: it quantifies \emph{why} a routing change
improves performance, separating information gain from mismatch reduction.

\paragraph{Test-time compute scaling.}
The discovery that repeated sampling can compensate for model
size~\citep{brown2024monkeys,snell2025testtime} motivates the
crossover formula~\eqref{eq:crossover}.
\citet{wu2025inference} derive inference scaling laws relating compute
to accuracy; \citet{osca2024} optimize sample-compute allocation.
Our crossover formula makes the tradeoff explicit at the level of individual
task modes, rather than in aggregate.
The key distinction is that these works study \emph{how much}
test-time compute to use; our decomposition studies \emph{which system
configuration} to use, with test-time compute as one of several levers.

\paragraph{Algorithm selection and transfer.}
The algorithm selection problem~\citep{rice1976algorithm} asks which
solver to apply to a given problem instance.
SATzilla~\citep{xu2012satzilla} and AutoFolio~\citep{lindauer2015autofolio}
solve this with portfolio-based methods; SMAC~\citep{hutter2011smac} and
Hyperband~\citep{li2017hyperband} use sequential model-based or
bandit-based optimization.  These treat each solver as a black box.
Our separable evaluator exploits the packed-family structure to score
solvers analytically from shared pilot data, achieving $O(1)$-in-$k$
evaluation cost where portfolio methods require $O(k)$.

Transfer of algorithm configurations across problem
distributions connects to domain adaptation
theory~\citep{bendavid2010theory,mansour2009domain} and
transferability estimation~\citep{achille2019task2vec,you2021logme,
pandy2022transferability}.  The Ben-David theory bounds generalization
error under distribution shift via $\mathcal{H}$-divergence; our transfer bounds
use the more structured $(\alpha_{\mathrm{PT}}, \alpha_O)$ divergence
pair, which exploits the decomposition to separate competence shift from
routing shift.  This finer-grained analysis enables the three-case
decision rule, which has no analogue in generic domain adaptation.

Racing algorithms~\citep{maron1993hoeffding,maron1997racing} and
successive halving~\citep{jamieson2016successive} reduce evaluation cost
by eliminating poor candidates early.  They remain black-box methods:
each candidate must be evaluated individually.
Algorithm~\ref{alg:decomp-guided} is white-box in the sense that it
exploits the algebraic structure of the performance objective to avoid
per-candidate evaluation entirely in Phase~2.
Weitzman's Pandora's box rule~\citep{weitzman1979optimal} solves optimal
sequential search with inspection costs; our setting differs in that the
separable evaluator makes all inspections simultaneous.

\section{Discussion}
\label{sec:discussion}

\paragraph{Limitations.}
The packed-family assumption is the central modeling commitment.
It requires that task difficulty clusters into discrete modes with
geometric time distributions---a reasonable approximation for
verifiable tasks with independent attempts, but a simplification for
tasks with complex dependencies between retries (e.g., where each
attempt builds on the previous one).  Outside the packed family, the
four-term identity holds only approximately, with a bounded residual
that depends on departure from the geometric-trials
model~\citep{beneventano2026framework}.

The transfer thresholds are conservative by design.  The 73.8\% accuracy
on borderline cells reflects worst-case Pinsker and entropy-continuity
bounds; tighter bounds (e.g., Bernstein-type refinements or empirical
validation on diagnostic samples) would likely improve accuracy.  The
thresholds should be treated as guidelines, not hard boundaries.

The framework is most valuable when models have comparative advantages across modes and $k \gg m \cdot \Mn$; when one model dominates everywhere (as in EXP15), the decomposition correctly diagnoses $G \approx 0$ but offers no advantage over simple model selection.

The design space in our experiments is finite and low-dimensional
(model $\times$ depth).  Extending to continuous design spaces---joint
optimization over model mixtures, context budgets, and tool
configurations---is an important open direction.

\paragraph{Takeaways.}
Three things are clearer after this work.
First, the separation of cost, competence, information, and routing
is not merely a diagnostic convenience---it is the algebraic structure
that makes $O(1)$-in-$k$ design selection possible.
Second, the two dimensions of transfer (competence and orchestration)
are genuinely independent, which means that a single similarity score
between task domains is insufficient for transfer decisions.
Third, even when a framework's assumptions are violated (as in EXP15,
where routing is trivial), the decomposition correctly diagnoses the
violation and remains useful for scoping.

\paragraph{Open questions.}
Can the sample-complexity cascade be tightened beyond the Bernstein
refinement?  Empirical pilot sizes are roughly $29\times$ smaller than
the Hoeffding bound predicts; closing this gap would make the budget
prescriptions practical rather than merely conservative.
Can the framework be extended to tasks where attempts are not independent
(e.g., agentic loops with memory)?  The geometric-trials assumption is
the binding constraint.
Finally, validating the decomposition and transfer theory on production
agent deployments---where mode structure must be discovered rather than
assumed---is the most important next step.

\paragraph{Acknowledgements.}
This work was conducted using the PoggioAI/MSc research pipeline
\citep{PoggioAIMSc_2026}, an experimental framework for ML theory
research with humans on the loop.


\bibliographystyle{unsrtnat}
\bibliography{references}

\newpage


% ==================================================================
%  APPENDIX (unlimited)
% ==================================================================

\appendix

\section{Proof of the Four-Term Identity (Theorem~\ref{thm:four-term})}
\label{app:four-term}

Under Assumption~\ref{ass:packed}, the expected log certified time for
design $(M, q)$ at depth $D$ is:
\begin{equation}
  \E[-\log T] = \sum_s \pi_s \log\bigl(1 - (1 - p(s,M))^D\bigr)
    + \sum_s \pi_s \log q(s).
\end{equation}
The first sum is the competence contribution; the second is the routing
contribution.  The log-performance gap between $(M_0, q_0)$ and $(M, q)$
is:
\begin{align}
  \Delta &= \E[-\log T(M,q)] - \E[-\log T(M_0, q_0)] \nonumber \\
  &= \underbrace{\log\frac{\kappa_0}{\kappa(M)}}_{\text{cost}}
    + \underbrace{\sum_s \pi_s \log\frac{p(s,M)}{p(s,M_0)}}_{\varphi}
    + \underbrace{\sum_s \pi_s \log\frac{q(s)}{q_0(s)}}_{\text{routing difference}}.
\end{align}
The routing difference decomposes via the chain rule for KL divergence.
Define the posterior mode distribution $\pi_D^M(s) \propto \pi_s \cdot
P_D(s,M)$ where $P_D(s,M) = 1 - (1-p(s,M))^D$.  Then:
\begin{align}
  \sum_s \pi_s \log\frac{q(s)}{q_0(s)}
  &= \MI_\pi(S; D_M)
    - \E_D\bigl[\KL(\pi_D^M \| q_D)\bigr]
    + \text{constant w.r.t.\ } q_0 \nonumber \\
  &= G - \mismatch + \text{const}.
\end{align}
Combining yields the four-term identity~\eqref{eq:four-term}. \qed

\section{Lagrangian Non-Connection}
\label{app:lagrangian}

Consider the packed-family optimization with $\Mn{=}2$ modes, two models,
and a cost budget~$B$.  The Lagrangian is:
\[
  \mathcal{L}(q, M, \lambda) = \E[-\log T(M, q)]
    + \lambda\bigl(\kappa(M) \cdot D - B\bigr).
\]
The KKT conditions yield a single scalar multiplier $\lambda^* \ge 0$
as the shadow price of the budget constraint.  The optimal routing satisfies
$q_s^* = \pi_s$ (posterior-matched), and the optimal objective equals
$\Hent(\pi)$, the mode entropy, independent of~$M$.

The four decomposition terms $(\text{cost}, \varphi, G, \mismatch)$ are
components of the primal objective evaluated at any feasible design
point---a zeroth-order algebraic property.  Shadow prices measure
sensitivity to constraint relaxation---a first-order envelope property.
No linear combination of the single shadow price $\lambda^*$ reproduces
the four-term decomposition, except in degenerate cases where the budget
constraint is not binding.

The structural reason is that the decomposition is a \emph{pointwise
identity} valid at every $(M, q)$, not a first-order optimality condition
valid only at the optimum.  The decomposition's value to optimization is
as a separable evaluator, not as a dual decomposition. \qed

\section{Near-Optimality Proof (Proposition~\ref{prop:near-opt})}
\label{app:near-opt}

\begin{proof}
Since $\hat{d}$ minimizes $\hat{f}$ over $\cD$:
\begin{align}
  f(\hat{d})
  &\le \hat{f}(\hat{d}) + \delta
    \tag{$\delta$-accuracy for $\hat{d}$} \\
  &\le \hat{f}(d_{\mathrm{opt}}) + \delta
    \tag{$\hat{d}$ is best estimated} \\
  &\le f(d_{\mathrm{opt}}) + 2\delta.
    \tag{$\delta$-accuracy for $d_{\mathrm{opt}}$}
\end{align}

\textbf{Tightness.}
Take $k{=}2$, $f(d_1) = 0$, $f(d_2) = 2\delta - \epsilon$.
Set $\hat{f}(d_1) = +\delta$, $\hat{f}(d_2) = \delta - \epsilon$.
Algorithm selects $d_2$; regret $= 2\delta - \epsilon \to 2\delta$.

\textbf{Confirmation bound.}
Since $\hat{d} \in \cC$, $\tilde{f}(d^*) \le \tilde{f}(\hat{d})$.
Applying $\eta$-accuracy:
$f(d^*) \le \tilde{f}(d^*) + \eta \le \tilde{f}(\hat{d}) + \eta
\le f(\hat{d}) + 2\eta \le f(d_{\mathrm{opt}}) + 2\delta + 2\eta$. \qed
\end{proof}

\section{Evaluation Complexity Proof}
\label{app:eval-complexity}

Phase~1: $m$ models $\times$ $\Mn$ modes $\times$ $n_0$ instances
$= m \cdot \Mn \cdot n_0$ evaluations.
Phase~2: $k \cdot \Mn$ arithmetic operations, zero evaluations.
Phase~3: $c \cdot N_{\mathrm{conf}}$ evaluations.
Total: $N_{\mathrm{Alg1}} = m \cdot \Mn \cdot n_0 + c \cdot N_{\mathrm{conf}}$.

Racing: $k \cdot n_{\mathrm{eval}}$ evaluations (with
$n_{\mathrm{eval}} = O(\log(k/\beta)/\delta^2)$ for Hoeffding-based
accuracy).

Savings factor:
$N_{\mathrm{race}}/N_{\mathrm{Alg1}} = k \cdot n_{\mathrm{eval}} /
(m \cdot \Mn \cdot n_0 + c \cdot N_{\mathrm{conf}})$.
For fixed $m, \Mn, c, N_{\mathrm{conf}}$, this is $\Theta(k)$.

\textbf{Asymptotic:}
$N_{\mathrm{Alg1}} = O(1)$ in $k$;
$N_{\mathrm{race}} = O(k)$;
$N_{\mathrm{SHA}} = O(B)$ with first-round per-design budget
$O(B/(k\log k)) \to 0$. \qed

\section{Sample Complexity Cascade Proof}
\label{app:cascade}

Chain three results.  From Hoeffding applied to $m \cdot \Mn$ mode--model
pairs with $n_0 = N/(m \cdot \Mn)$ samples each:
\[
  \delta \le \sqrt{\frac{\log(2\Mn/\dconf)}{2n_0}}
  = \sqrt{\frac{m \cdot \Mn \cdot \log(2\Mn/\dconf)}{2N}}.
\]
From Proposition~\ref{prop:near-opt}: regret $\le 2\delta$.
Substituting:
\[
  \text{regret} \le 2\sqrt{\frac{m \cdot \Mn \cdot \log(2\Mn/\dconf)}{2N}}.
\]
Bernstein refinement: replace $1/(2n_0)$ with
$3\pmin/n_0$ when $\pmin \le 1/2$:
\[
  \text{regret}_{\mathrm{Bern}} \le
  2\sqrt{\frac{3m \cdot \Mn \cdot \pmin \cdot \log(2\Mn/\dconf)}{N}}.
\]
Improvement factor: $\sqrt{6\pmin}$; at $\pmin = 0.02$, this is
$\approx 0.35$, a $2.9\times$ tighter bound. \qed

\section{Independence Proofs (Proposition~\ref{prop:independence})}
\label{app:independence}

\subsection{C64: Post-training equivalence does not imply orchestration equivalence}

\textbf{Construction.}
$\Mn{=}2$, $m{=}2$ ($M_{\mathrm{big}}, M_{\mathrm{small}}$).
Shared competence:
$p(1, M_{\mathrm{big}}) = 0.8$, $p(2, M_{\mathrm{big}}) = 0.3$,
$p(1, M_{\mathrm{small}}) = 0.4$, $p(2, M_{\mathrm{small}}) = 0.6$.
Mode distributions: $\pi_A = (0.7, 0.3)$, $\pi_B = (0.3, 0.7)$.

Post-training equivalence holds (identical per-mode, per-model success
probabilities).

Orchestration equivalence fails: $\KL(\pi_A \| \pi_B) = 0.339 > 0$.
The posterior routing distributions at $D{=}1$ are
$\pi_1^A \approx (0.757, 0.243)$ vs.\
$\pi_1^B \approx (0.364, 0.636)$,
with $\KL(\pi_1^A \| \pi_1^B) \approx 0.320$.
The ordering persists at all depths by monotonicity of
$P_D(s) = 1 - (1-p^*(s))^D$.

\subsection{C65: Orchestration equivalence does not imply post-training equivalence}

\textbf{Construction.}
$\pi_C = \pi_D = (0.5, 0.5)$.
$\tau_C$: $p_C(1, M_{\mathrm{big}}) = 0.9$,
  $p_C(2, M_{\mathrm{big}}) = 0.2$,
  $p_C(1, M_{\mathrm{small}}) = 0.3$,
  $p_C(2, M_{\mathrm{small}}) = 0.7$.
$\tau_D$: $p_D(1, M_{\mathrm{big}}) = 0.6$,
  $p_D(2, M_{\mathrm{big}}) = 0.1$,
  $p_D(1, M_{\mathrm{small}}) = 0.2$,
  $p_D(2, M_{\mathrm{small}}) = 0.8$.

Orchestration equivalence holds at the routing-policy level:
$M_{\mathrm{big}}$ wins mode~1 and $M_{\mathrm{small}}$ wins mode~2
in both distributions, at all depths (by strict monotonicity).

Post-training equivalence fails:
$|p_C(1, M_{\mathrm{big}}) - p_D(1, M_{\mathrm{big}})| = 0.3$.
The competence term $\varphi$ differs by $|\varphi_C - \varphi_D|
\approx 0.413$.

\section{Three-Case Transfer Proof (Theorem~\ref{thm:transfer})}
\label{app:transfer}

\textbf{Setup.}
Write the transfer error as:
\[
  \Delta = f(d_S^*; \tau_T) - f(d_T^*; \tau_T)
  \le |f(d_S^*; \tau_T) - f(d_S^*; \tau_S)|
    + |f(d_S^*; \tau_S) - f(d_T^*; \tau_T)|.
\]
Since $d_S^*$ optimizes $\tau_S$: $f(d_S^*; \tau_S) \le f(d_T^*; \tau_S)$.
It suffices to bound $|f(d; \tau_T) - f(d; \tau_S)|$ term by term.

\textbf{Cost term:} $|C_T - C_S| = 0$ (model-dependent, not
distribution-dependent).

\textbf{Competence term:}
$|\varphi_T - \varphi_S| \le 2\alpha_{\mathrm{PT}}$ (each of two
log-terms shifts by at most $\alpha_{\mathrm{PT}}$).

\textbf{Information gain:}
$|G_T - G_S| \le \sqrt{2\alpha_O} \cdot \log\Mn$ by Pinsker's inequality
($\|\pi_S - \pi_T\|_{\mathrm{TV}} \le \sqrt{\alpha_O/2}$) and the
Audenaert entropy-continuity bound.

\textbf{Routing mismatch:}
$|\mismatch_T - \mismatch_S| \le \sqrt{2\alpha_O} \cdot \log\Mn$
by the same Pinsker argument applied to the KL divergence of routing
allocations.

\textbf{Case~1:}
$\Delta \le 2\alpha_{\mathrm{PT}} + \sqrt{2\alpha_O}\cdot\log\Mn$.
Under $\alpha_{\mathrm{PT}} \le \tau_1 = \dtarget/2$ and
$\alpha_O \le \tau_2 = \dtarget^2/(2\log^2\Mn)$:
$\Delta \le \dtarget + \dtarget = 2\dtarget$.

\textbf{Case~2:}
Keep $M_S^*$, re-estimate routing from target pilot.
Residual error: $\alpha_{\mathrm{PT}}$ (competence) $+ 2\dacc$
(re-estimation, by Proposition~\ref{prop:near-opt} applied to the
restricted design space $\{(M_S^*, D) : D \in D_{\mathrm{values}}\}$).
Cost: $\Mn \cdot n_0'$ for one model instead of $m \cdot \Mn \cdot n_0$
for all models---a factor of $m$ savings.

\textbf{Case~3:}
No transfer. Full Algorithm~\ref{alg:decomp-guided} on $\tau_T$.
$\Delta \le 2\dacc$ by Proposition~\ref{prop:near-opt}. \qed

\section{Decision Threshold Derivation}
\label{app:thresholds}

Set each term in the Case~1 bound equal to $\dtarget$:
\begin{align}
  2\alpha_{\mathrm{PT}} &\le \dtarget
    &\implies \tau_1 &= \dtarget/2, \\
  \sqrt{2\alpha_O}\cdot\log\Mn &\le \dtarget
    &\implies \tau_2 &= \dtarget^2/(2\log^2\Mn).
\end{align}
For $\dtarget = 0.1$ and $\Mn = 5$:
$\tau_1 = 0.05$, $\tau_2 \approx 0.0019$.
The orchestration threshold $\tau_2$ is much smaller than $\tau_1$,
reflecting the amplification by $\log\Mn$: even small mode-distribution
shifts can cause significant routing errors.

\section{Transfer Savings Proof (Proposition~\ref{prop:transfer-savings})}
\label{app:transfer-savings}

Case~1: $\mathcal{C}_{\mathrm{target}} = 0$.
Case~2: Pilot evaluates one model on $\Mn$ modes instead of $m$ models.
Pilot cost: $\Mn \cdot n_0$ vs.\ $m \cdot \Mn \cdot n_0$.
Ratio: $1/m$.
Case~3: $\mathcal{C}_{\mathrm{target}} = m \cdot \Mn \cdot n_0 +
c \cdot N_{\mathrm{conf}} = \mathcal{C}_S$. No savings. \qed

\section{Additional Experiment Details}
\label{app:experiments}

\subsection{EXP10 Setup}

Five models (tiny, small, medium, large, xlarge) with costs
$\kappa \in \{0.5, 1, 3, 10, 25\}$.
Per-mode success probabilities scale linearly with model size.
$\Mn{=}5$ modes with Zipf distribution
$\pi \propto (1, 1/2, 1/3, 1/4, 1/5)$.
Cost--quality weight $\mu = 0.3$.
Pilot: $n_0 = 200$ per mode per model.
Confirmation: $c = 5$, $N_{\mathrm{conf}} = 500$.
Racing: $n_{\mathrm{eval}} = 500$ per design.
50 Monte Carlo replications per condition.

\subsection{EXP13 Setup}

Same model and mode parameters as EXP10.
All methods given a matched budget of 7,500 total evaluations.
BO: RBF kernel, length scale 0.5, EI acquisition, 5 initial random
points, 150 evaluations per query.
SHA: budget $= 7{,}500$, $\lceil\log_2 k\rceil$ rounds.
50 MC replications per condition.

\subsection{EXP14 Setup}

Three models (specialist\_A, generalist\_B, specialist\_C) with
heterogeneous competence profiles.
$k = 51$ designs (3 models $\times$ 17 depths).
$\dtarget = 0.15$, yielding $\tau_1 = 0.075$, $\tau_2 \approx 0.0043$.
100 MC replications per grid cell.

\subsection{EXP15 Data Sources}

Semi-synthetic pass rates calibrated to published SWE-bench-lite
results~\citep{jimenez2024swebench}.
Repository-specific pass rates derived from leaderboard entries
(circa early 2025).
Not exact per-instance predictions; structural findings (mode
heterogeneity, absence of crossing structure) are robust to the
exact numbers.


\end{document}
